% TOP_PROBLEM: {"problemId": "problem.p-versus-bpp-does-bpp-p", "canonicalTitle": "P versus BPP (Does BPP = P?)", "releaseRank": 27, "primaryDomain": "cryptography_coding_information_optimization", "primaryDomainLabel": "Cryptography, coding, information & optimization", "displayStatus": "open", "releaseStatus": "open"}
% !TeX program = lualatex
% ATTEMPT_STATUS: unresolved
\documentclass[11pt]{article}
\usepackage[margin=25mm]{geometry}
\usepackage{fontspec}
\setmainfont{Latin Modern Roman}
\usepackage{amsmath,amssymb,mathtools}
\usepackage{unicode-math}
\setmathfont{Latin Modern Math}
\usepackage{xurl,hyperref}
\hypersetup{colorlinks=true,urlcolor=blue}
\setcounter{secnumdepth}{0}
\setlength{\parindent}{0pt}
\setlength{\parskip}{5pt}
\setlength{\emergencystretch}{3em}
\begin{document}
\section*{\texorpdfstring{P versus BPP (Does BPP = P?)}{P versus BPP (Does BPP = P?)}}\label{27}
\subsection{Definitions and mathematical statement}
A language $L$ is in $\mathsf{BPP}$ if a randomized polynomial-time algorithm $A$ satisfies $\Pr[A(x)=1_L(x)]\ge2/3$ for every input $x$. Membership in $\mathsf P$ requires a deterministic polynomial-time decision algorithm with no errors. Determine whether
\[
 \forall L\in\mathsf{BPP},\quad L\in\mathsf P.
\]
The simulation must work in the ordinary uniform model, not merely through nonuniform advice or for most inputs.
\par\smallskip\textit{Formulation and concept references:} \href{https://people.seas.harvard.edu/~salil/pseudorandomness/pseudorandomness-Aug12.pdf}{[S1]}.

\subsection{Short English statement}
Can randomness always be removed from an efficient bounded-error decision algorithm without losing polynomial-time efficiency?
\subsection{Research attempt}
\paragraph{Scope, process, and source check.}
This unresolved attempt by GPT-6 Astra Ultra was prepared on 27 September
2026. The method was to isolate the randomness actually needed by a fixed
bounded-error decider, derive sufficient replacement conditions, and test
whether existential or limited-independence arguments meet those conditions.
All inputs of a given length must be handled, and the deterministic simulator
must be one finite uniform algorithm. Neither advice depending on the length
nor an average-case guarantee is substituted for that requirement.

The pseudorandomness framework in Vadhan [S1] distinguishes short seed length,
fooling guarantees, and efficient construction. Adleman's classical
nonuniform simulation [E1] supplies the relevant warning: bounded-error
randomness can be fixed separately at each input length without giving a
uniform derandomization. Impagliazzo--Wigderson [E2], Theorem 2, prove
$\mathsf P=\mathsf{BPP}$ under the hypothesis that some language in
$\mathsf E=\mathsf{DTIME}(2^{O(n)})$ has worst-case circuit complexity
$2^{\Omega(n)}$. That hypothesis is not proved here. The catalogue's
recent result [S3] treats insensitive algorithms; no unrestricted simulation
is inferred from its restricted algorithm class.

\paragraph{Attack plan and normalization.}
Fix a BPP decider $A$ for a total language $L$. On inputs of length $n$, pad
its random tape to a deterministic polynomial length $m(n)$; unused bits
have no effect. Write $A(x;r)\in\{0,1\}$ for the resulting deterministic
computation. Its time is polynomial for every $r$, not merely in expectation.
A sufficiently good, uniformly generated polynomial-size collection of
random tapes would permit deterministic enumeration. The main route is to
state exactly what ``sufficiently good'' means, then compare that condition
with what amplification and independence really provide.

The containment $\mathsf P\subseteq\mathsf{BPP}$ follows by ignoring random
bits. Thus an upper-bound solution needs the reverse containment for every
fixed BPP decider, with a polynomial runtime allowed to depend on that decider.
A negative solution needs one total bounded-error language outside every
uniform deterministic polynomial-time bound.

\paragraph{Lemma 27.1: a uniform short-seed condition is sufficient.}
Suppose that for the fixed $A$ there is a deterministic algorithm computing
\[
 G_n:\{0,1\}^{d(n)}\longrightarrow\{0,1\}^{m(n)},
 \qquad d(n)\le c\log_2(n+2),
\]
in time polynomial in $n$, where $d(n)$ is efficiently computable and $c$
is a constant. Suppose also that, for every $x\in\{0,1\}^n$,
\begin{equation}\label{a27:fooling}
 \left|\Pr_{s}[A(x;G_n(s))=1]-\Pr_r[A(x;r)=1]\right|\le1/12.
\end{equation}
Then $L\in\mathsf P$.

\emph{Proof.} Enumerate all $2^{d(n)}\le(n+2)^c$ seeds, compute their
outputs, and count accepting simulations. On yes-inputs the fraction is
at least $2/3-1/12=7/12$. On no-inputs it is at most
$1/3+1/12=5/12$. Comparing the integer acceptance count with half the number
of seeds therefore decides $L$ without error. The generator and each
simulation take polynomial time, as does the number of seeds. Small input
lengths are included by the displayed $(n+2)$ bound, or can be hardwired
separately if the hypothesis begins only eventually.\hfill$\square$

This lemma asks only to fool the tests $r\mapsto A(x;r)$ associated with one
decider, not every Boolean function on the random tape. The generator may
depend on $A$, but its algorithm is fixed before seeing $x$. It is the
uniformity and simultaneous guarantee for all $x$ that will be missing below.
A generator with seed length $(\log n)^2$ would instead give the enumeration
cost $2^{O((\log n)^2)}$, which is not a polynomial bound. A short seed is
also insufficient if evaluating one generator output takes superpolynomial
time. These costs cannot be hidden inside the name ``pseudorandom.''

\paragraph{Lemma 27.2: fixing randomness does give polynomial advice.}
For each length $n$ there is a list of
\[
                       R=12(n+2)+1
\]
random tapes such that majority vote among the corresponding executions of
$A$ is correct on every length-$n$ input.

\emph{Proof.} Select the $R$ tapes independently and uniformly. Fix an input
$x$ and let $X_i$ indicate an incorrect answer on tape $i$. The variables
are independent and satisfy $\Pr[X_i=1]\le1/3$. Therefore
\[
 \mathbb E\,2^{\sum_i X_i}
       =\prod_i\mathbb E\,2^{X_i}\le(4/3)^R.
\]
Markov's inequality gives
\[
 \Pr\left[\sum_iX_i\ge R/2\right]
       \le\left(\frac4{3\sqrt2}\right)^R
       <(8/9)^{6(n+2)}<2^{-(n+2)}.
\]
The last inequality uses $(8/9)^6<1/2$, or equivalently
$2\cdot8^6<9^6$. As $R$ is odd, majority is unambiguous. There are $2^n$
possible inputs, so the union bound makes the probability of any wrong
majority less than $1/4$. At least one list has no wrong majority at all.
Its length is $Rm(n)$, a polynomial in $n$.\hfill$\square$

Hardwire such a list at each length and simulate the fixed algorithm by
polynomial-size circuits. This proves the familiar containment
$\mathsf{BPP}\subseteq\mathsf{P/poly}$. The displayed estimate was derived
explicitly to make its exact scope visible. The list is independent of the
particular input, but its existence proof does not tell a uniform algorithm
which list to print. Brute-force enumeration of candidate lists followed by
checking all $2^n$ inputs is not a polynomial-time construction; checking
correctness also requires knowing the correct language answers. No efficient
replacement for that verification is supplied here.

The two lemmas consequently do not combine automatically. Lemma 27.1 has
an effective generator hypothesis, while Lemma 27.2 provides only existential
length-dependent data. Treating the latter as a subroutine would introduce
precisely the nonuniform advice forbidden in the target statement.

\paragraph{An exact test of the limited-independence repair.}
One might try to construct replacement tapes by ensuring that all small sets
of random bits are independent. The following example shows that even a
very strong version of that property does not fool every polynomial-time
algorithm.

For $m\ge2$, let $D_m$ be uniform over pairs $(u,v)\in\{0,1\}^m\times
\{0,1\}^m$ such that both blocks have even parity. It is efficiently sampled:
choose the first $m-1$ bits of each block freely and set the last bit to
their XOR. Every set of at most $m-1$ coordinates in the full $2m$-bit string
is uniformly distributed. Indeed, in each block any prescribed values on
fewer than $m$ coordinates have exactly $2^{m-t-1}$ even-parity completions,
where $t$ coordinates were prescribed. The two blocks are independent,
so the joint marginal is also uniform.

Now consider the algorithm that accepts exactly when both blocks have even
parity. With fully independent fair bits its acceptance probability is
$1/4$. It is therefore a valid BPP decider for the empty language, with
success $3/4$ on every input; choose $m$ polynomial in input length and
ignore the input contents. Under $D_m$ it accepts with probability $1$ and
is wrong on every input. Repeating executions with independently drawn
$D_m$ tapes preserves this certain error.

This is a counterexample to a proposed universal derandomization rule,
not a hard language: the empty language has the trivial deterministic
rejecting algorithm. It shows that many-wise independence constrains local
marginals while a linear-time parity test can inspect a global correlation.
The proof of a pseudorandom generator must control the actual computational
tests in \eqref{a27:fooling}; replacing that proof by an independence count
is invalid. Restricted algorithms that cannot detect such correlations may
admit different positive results, with their hypotheses retained.

\paragraph{Why an input-dependent seed search is not yet an algorithm.}
There is a second quantifier trap in asking for one fixed logarithmic-seed
generator to fool \emph{all} polynomial-time tests at once. Suppose its
output length is $m(n)=n$ and its range has at most $(n+2)^c$ strings.
The test ``is $r$ in the range of $G_n$?'' enumerates every seed and compares
its generated output with $r$. If the generator is polynomial-time, this
test is also polynomial-time. It accepts generator outputs with probability
$1$, but accepts a uniform $n$-bit string with probability at most
$(n+2)^c/2^n$, eventually less than $1/3$. On those lengths it is itself a
valid randomized decider for the empty language; finitely many smaller
lengths can simply reject. Thus that fixed generator fails this particular
efficient test.

This does not obstruct Lemma 27.1: its generator may depend on the decider
or on a specified circuit-size bound. Increasing the test's polynomial
exponent may require increasing the seed constant or the generator's work.
The support-membership test proves that those dependencies cannot be erased
by asserting one fixed polynomial budget that works for every polynomial
budget. It is also another counterexample to a proposed universal replacement
rule, not evidence that the empty language needs randomness.

For a yes-input, many tapes lead to acceptance; for a no-input, many lead to
rejection. Finding an accepting tape alone therefore does not identify the
answer. A BPP machine may have accepting tapes on both sorts of input.
Likewise, finding a tape on which repeated runs agree is insufficient unless
that agreement is certified to reflect the true majority over all tapes.
The two-sided error convention must remain in the argument.

A useful exact restricted case is $m(n)=O(\log n)$: enumerate every original
random tape and compare the true acceptance fraction with $1/2$. No generator
is needed, and the result is a uniform polynomial-time decider. The missing
step for arbitrary BPP is thus not amplification alone. It is a method of
replacing polynomially many random bits by logarithmically many effectively
enumerable bits while preserving the decision gap, or another simulation
achieving the same effect.

One cannot choose a successful tape list for each input by an unrestricted
existential quantifier and then call that choice a deterministic computation.
Nor can the list be fixed once for all input lengths: its tape lengths and
coverage obligations change with $n$. Lemma 27.2 addresses each finite input
length separately, as nonuniform circuit complexity allows.

\paragraph{Verification, remaining gap, and outcome.}
The companion script checks the amplification bound using exact rational
binomial probabilities at the boundary failure probability $1/3$, verifies
$2\cdot8^6<9^6$, and tests all coordinate marginals of $D_m$ through $m=5$.
It confirms acceptance $1/4$ under uniform tapes and $1$ under the replacement.
These are finite checks of the displayed claims; their all-length versions
follow from the proofs. No heuristic randomness test is treated as a proof
that a generator fools all polynomial-time computations.

\textbf{UNRESOLVED.} The strongest uniform conclusion proved here is exact
derandomization for logarithmic randomness, and the conditional short-seed
lemma. The all-input tape-fixing argument remains nonuniform, and the proposed
limited-independence replacement has an explicit failing decider. The first
remaining gap is a uniform polynomial-time construction meeting the necessary
all-input gap guarantee for an arbitrary BPP decider. A lower-bound solution
would instead need one BPP language outside every deterministic polynomial
bound. Neither is established, and no novelty is claimed for the elementary
or classical calculations.

\subsection{Sources}
\begin{itemize}
\item[S1]S. Vadhan, 'Pseudorandomness', Foundations and Trends in Theoretical Computer Science 7(1-3):1-336, 2012.
\url{https://people.seas.harvard.edu/~salil/pseudorandomness/pseudorandomness-Aug12.pdf}.
\textit{Formulation source.}
\item[S2]CS 6810: Theory of Computing, Lecture 15, March 17, 2026.
\url{https://www.cs.cornell.edu/courses/cs6810/2026sp/lec15.pdf}.
\textit{Further reference; not a proof endorsement.}
\item[S3]Pseudorandomness Beating the Hybrid Argument for Insensitive Algorithms.
\url{https://eccc.weizmann.ac.il/report/2026/082/}.
\textit{Further reference; not a proof endorsement.}
\item[E1]Leonard M. Adleman, \emph{Two Theorems on Random Polynomial Time},
Proceedings of FOCS 1978, pp. 75--83.
\url{https://doi.org/10.1109/SFCS.1978.37}.
\url{https://mikael-monet.net/share/papers/adleman1978.pdf}.
\textit{Primary nonuniform-simulation reference; consulted 27 September 2026.
The explicit amplification and tape-fixing calculation is given above.}
\item[E2]Russell Impagliazzo and Avi Wigderson,
\emph{P = BPP if E Requires Exponential Circuits: Derandomizing the XOR Lemma},
Proceedings of STOC 1997, pp. 220--229, Theorem 2.
\url{https://www.math.ias.edu/~avi/PUBLICATIONS/MYPAPERS/IW97/proc.pdf}.
\url{https://doi.org/10.1145/258533.258590}.
\textit{Conditional hardness-versus-randomness result; consulted
27 September 2026. Its circuit-hardness assumption is not asserted here.}
\end{itemize}
\end{document}
